\documentclass[10pt,a4paper]{amsart}
\usepackage[margin=19mm]{geometry}
\usepackage{amsmath,amssymb,mathtools}
\usepackage[scr=rsfs,cal=euler]{mathalfa}
\usepackage{xfrac}
\usepackage[unicode,hidelinks,bookmarks=false]{hyperref}
\newcommand{\ie}{i.e.\ }
\newcommand{\cf}{cf.\ }
\newcommand{\resp}{respectively\ }
\newcommand{\Subsection}{\subsection}
\providecommand{\nobreakdash}{\nobreak\hbox{-}}
\setlength{\emergencystretch}{2em}
\multlinegap0pt
\usepackage{bm}
\usepackage{bbm}
\usepackage{dsfont}
\usepackage{xspace}
\usepackage{cancel}
\usepackage{mathtools}
\usepackage[shortlabels]{enumitem}
\usepackage{amsthm}
\makeatletter
\@ifundefined{corollary}{\newtheorem{corollary}{Corollary}}{}
\@ifundefined{lemma}{\newtheorem{lemma}{Lemma}}{}
\@ifundefined{proposition}{\newtheorem{proposition}{Proposition}}{}
\@ifundefined{theorem}{\newtheorem{theorem}{Theorem}}{}
\makeatother
\DeclareMathOperator{\HH}{H}
\DeclareMathOperator{\Ob}{Ob}
\DeclareMathOperator{\reg}{reg}
\newcommand{\A}{{\mathscr{A}}}
\newcommand{\B}{{\mathscr{B}}}
\newcommand{\C}{{\mathscr{C}}}
\newcommand{\FI}{{\mathrm{FI}}}
\newcommand{\N}{{\mathbb{N}}}
\newcommand{\hor}{{\mathrm{hor}}}
\newcommand{\ver}{{\mathrm{ver}}}
\title[Bounding regularity of $FI^m$-modules]{Bounding regularity of $FI^m$-modules}
\author{Wee Liang Gan, Khoa Ta}
\date{Source: arXiv:2407.20428v2 (CC BY 4.0)}
\begin{document}
\maketitle

\noindent\emph{Source and licence.} Bounding regularity of $FI^m$-modules. Version arXiv:2407.20428v2 (CC BY 4.0), distributed under the Creative Commons Attribution 4.0 International licence, as recorded in the arXiv licence metadata for this exact version and shown on the abstract page https://arxiv.org/abs/2407.20428v2. Full rights notice: licenses/arxiv-2407.20428-CC-BY-4.0.txt.\par\smallskip

\noindent\textit{Editorial scope.} Counted unit: the Proposition whose source label is \texttt{prop:2} in the article (source0.tex lines 732--740, source label \texttt{prop:2}), delivered together with the article's own complete original proof of it (source0.tex lines 742--804, one \texttt{proof} environment of 63 lines). No mathematics is written, repaired or re-derived: statement and proof are the article's own text line for line. Required context retained uncounted: homology\_and\_restrict (lines 304--312); grothendieck (lines 325--331); cor:finite degree case (lines 643--646). Every definition, display and label of those lines is retained unchanged. Omitted from this package and not redistributed: 1--303, 313--324, 332--642, 647--731, 741--741, 805--1220. Nothing inside a retained excerpt is omitted or altered: every retained body is delivered exactly as the article's lines are written, so no omission receipt of any kind accompanies this package. Citations: the retained excerpts carry no citation command at all, so no bibliography entry of the article is transcribed and no reference is redistributed. Reference adaptation: none was needed and none was made; every reference inside the delivered unit resolves inside it, so no unresolved reference or citation is produced. Printed slips kept literal: no printed word, symbol, hypothesis or number of the counted statement or of its proof was altered. Layout and class: the article's own class is \texttt{amsart} with packages including amssymb, amsthm, amsmath, amsrefs, graphics, hyperref, xcolor, xy, eucal, enumerate, thmtools; the wrapper is the lane's pinned \texttt{amsart} editorial wrapper at 10pt on A4, which restates the theorem-like environments the retained text needs and adds the article's own macro definitions that the retained text needs (\texttt{\textbackslash A}, \texttt{\textbackslash B}, \texttt{\textbackslash C}, \texttt{\textbackslash FI}, \texttt{\textbackslash HH}, \texttt{\textbackslash N}, \texttt{\textbackslash Ob}, \texttt{\textbackslash hor}, \texttt{\textbackslash reg}, \texttt{\textbackslash ver}), with extra wrapper packages (bm, bbm, dsfont, xspace, cancel, mathtools, enumitem). The delivered numbering is the unit's own. Assets: no retained span contains a figure environment, a graphics-inclusion command, a PDF-page inclusion, a table or a TikZ picture; no image file is copied into this package, the package declares no asset dependency and no image file is redistributed with it. Licence: version arXiv:2407.20428v2 of this article is distributed under the Creative Commons Attribution 4.0 International licence, as recorded in the arXiv licence metadata for this exact version and shown on the abstract page https://arxiv.org/abs/2407.20428v2. A full rights notice is delivered at licenses/arxiv-2407.20428-CC-BY-4.0.txt. Characters: every retained excerpt is pure ASCII, so the delivered engine, XeLaTeX with its default Latin Modern text font, reports no missing character.
\begin{lemma} \label{homology_and_restrict}
   Let $V$ be a $\C$-module. 
   \begin{enumerate}[(i)]
      \item Let $X\in \Ob(\A)$. For each $i\geqslant 0$, we have
      \[ (\HH^{\ver}_i(V))_{(X,-)} \cong \HH^{\B}_i(V_{(X,-)}). \] 
      \item Let $Y\in \Ob(\B)$. For each $i\geqslant 0$, we have
      \[ (\HH^{\hor}_i(V))_{(-,Y)} \cong \HH^{\A}_i(V_{(-,Y)}). \] 
   \end{enumerate}
\end{lemma}

\begin{theorem} \label{grothendieck}
Let $V$ be a $\C$-module. Then there are two convergent first-quadrant spectral sequences:
\begin{gather*}
^{I}\!E^2_{pq} = \HH^\ver_p (\HH^\hor_q (V)) \Rightarrow \HH^{\C}_{p+q}(V),\\
^{II}\!E^2_{pq} = \HH^\hor_p (\HH^\ver_q (V)) \Rightarrow \HH^{\C}_{p+q}(V).
\end{gather*}
\end{theorem}

\begin{corollary} \label{cor:finite degree case}
   Let $V$ be an $\FI^m$-module. Assume that $\deg(V) < \infty$. Then 
   \[ \reg(V) \leqslant \deg(V). \] 
\end{corollary}

\begin{proposition} \label{prop:2}
   Let $m, \alpha, \gamma \in \mathbb{Z}$ with $m\geqslant 2$ and $\alpha, \gamma \geqslant -1$. Let $V$ be an $\FI^m$-module such that for all $x\in \N$:
   \begin{align*}
      x + \reg\left(V_{(x, -)}\right) &\leqslant \gamma \quad \mbox{ if } x\leqslant \alpha;\\
      V_{(x, -)} &= 0 \quad \mbox{ if } x > \alpha. 
   \end{align*}
   Then 
   \[ \reg(V) \leqslant \gamma. \]
\end{proposition}

\begin{proof}
   Applying Theorem \ref{grothendieck} to $V$ as a module over $\FI\times \FI^{m-1}$, we have the spectral sequence
   \[  ^{II}\!E^2_{pq} = \HH^\hor_p (\HH^\ver_q (V)) \Rightarrow \HH^{\FI^m}_{p+q}(V).  \]
   Let $\mathbf{n}=(n_1, n_2, \ldots, n_m) \in \N^m$ with 
   \[ n_1 + n_2 + \cdots + n_m > p+q+\gamma. \]
   It suffices to prove that 
   \[ (\HH^\hor_p (\HH^\ver_q (V)))_{\mathbf{n}} = 0. \]
   To this end, we first prove the following claim.
   
   \medskip 
   
   {\bf Claim \ref{prop:2}.1.} We have:
   \[ \deg\left( (\HH^\ver_q(V))_{(-, n_2, \ldots, n_m)} \right) < \max\{0, n_1-p\}. \]
   
   \begin{proof}[Proof of Claim \ref{prop:2}.1]
      Let $x\in \N$ with $x\geqslant n_1-p$. We need to show that 
      \[ (\HH^\ver_q(V))_{(x, n_2, \ldots, n_m)} = 0. \] 
      By Lemma \ref{homology_and_restrict}, we have
      \[ (\HH^\ver_q(V))_{(x, n_2, \ldots, n_m)} \cong (\HH^{\FI^{m-1}}_q(V_{(x, -)}))_{(n_2, \ldots, n_m)}. \]
      There are two cases: (1) $x\leqslant \alpha$, (2) $x>\alpha$.

      \medskip 


      {\bf Case 1: $x\leqslant \alpha$.}

      In this case we have:
      \begin{align*}
         n_2 + \cdots + n_m &> p+q+\gamma-n_1  \\
         &\geqslant p+q+x+  \reg\left(V_{(x, -)}\right) -n_1  \\
         &\geqslant p+q+(n_1-p) + \left(t_q\left(V_{(x,-)}\right) - q\right)-n_1 \\
         &\geqslant t_q\left(V_{(x,-)}\right).
      \end{align*}
      This implies that 
      \[ (\HH^{\FI^{m-1}}_q(V_{(x, -)}))_{(n_2, \ldots, n_m)} = 0. \]

      \medskip 



      {\bf Case 2: $x>\alpha$.} 

      In this case we have $V_{(x, -)}=0$, hence 
      \[ (\HH^{\FI^{m-1}}_q(V_{(x, -)}))_{(n_2, \ldots, n_m)} = 0. \]

      \medskip 

      This completes the proof of Claim \ref{prop:2}.1.
   \end{proof}

     Now by Lemma \ref{homology_and_restrict}, we have
      \[  (\HH^\hor_p (\HH^\ver_q (V)))_{\mathbf{n}} \cong (\HH^{\FI}_p( (\HH^\ver_q (V) )_{(-, n_2, \ldots, n_m)}  ))_{n_1}. \]
   
      If $n_1 \leqslant p$, then Claim \ref{prop:2}.1 implies that 
      \[ (\HH^\ver_q(V))_{(-, n_2, \ldots, n_m)} = 0, \] 
      so 
      \[ (\HH^{\FI}_p( (\HH^\ver_q (V) )_{(-, n_2, \ldots, n_m)}  ))_{n_1} = 0. \]

      If $n_1 > p$, then Claim \ref{prop:2}.1 implies that 
      \[ n_1 > p + \deg\left( (\HH^\ver_q(V))_{(-, n_2, \ldots, n_m)} \right). \]
      Hence by Corollary \ref{cor:finite degree case}, we have 
      \[ (\HH^{\FI}_p((\HH^\ver_q(V))_{(-, n_2, \ldots, n_m)}))_{n_1} = 0. \] 
\end{proof}

\end{document}
